\documentclass[mathserif,usenames,dvipsnames]{beamer}

%\usetheme{LCOM}
\usetheme{Madrid}
\usepackage{hyperref}
\usepackage{listings}
\usepackage{graphicx}
\usepackage[portuguese]{babel}
\usepackage[utf8]{inputenc}
\usepackage[T1]{fontenc}
\usepackage{booktabs}
\usepackage{amsmath,amssymb,amsthm,mathrsfs,amsfonts,dsfont}
\usepackage{adjustbox}
\usepackage{array}
\usepackage{longtable}
\usepackage{psfrag}
\usepackage{caption}
\usepackage{placeins}
\usepackage{subfig}
\usepackage{multirow}
\usepackage[siunitx, american]{circuitikz}

\uselanguage{portuguese}
\languagepath{portuguese}
\deftranslation[to=portuguese]{Theorem}{Teorema}
\deftranslation[to=portuguese]{theorem}{teorema}

\definecolor{green_variavel}{rgb}{0.13,0.57,0.2}
\usecolortheme[named=green_variavel]{structure}

%\definecolor{red_variavel}{rgb}{0.5,0,0}
%\usecolortheme[named=red_variavel]{structure}

\newcommand{\wait}{\vfill}
%\newcommand{\wait}{\pause}
	
\DeclareMathAlphabet\mathbfcal{OMS}{cmsy}{b}{n}
\newcommand\blfootnote[1]{%
	\begingroup
	\renewcommand\thefootnote{}\footnote{#1}%
	\addtocounter{footnote}{-1}%
	\endgroup
}

\setbeamertemplate{frametitle continuation}[from second][]

\title[Funções Singulares]{Funções Singulares e suas aplicações} 
\author[Vinícius Lagrota]{Vinícius Lagrota Rodrigues da Costa}
\institute[CES]{\includegraphics[width=0.3\textwidth]{figuras/ces}\\[\medskipamount]
	Centro de Ensino Superior de Juiz de Fora}
\date{\today}

\begin{document}

\begin{frame}
\maketitle 
\end{frame}

\begin{frame}
\frametitle{Sumário}
\tableofcontents
\end{frame}

%------------------------------------------------------------------------------------
\section{Recapitulação}
%------------------------------------------------------------------------------------
\begin{frame}
\frametitle{Recapitulação}
\framesubtitle{Resistores, capacitores e indutores}
\begin{minipage}[b]{0.3\linewidth}
	\begin{block}{Resistores}
		\begin{center}
			\begin{circuitikz}[european voltages]
				%				\draw (0,0) to [I, l^=$i(t)$] (0,2);
				\draw [-latex] (0.6, 1.85) -- (1.25, 1.85);
				\draw node[] at (0.25, 1.75) {$i(t)$};
				\draw (0,2) -- (1.5,2);
				\draw (0,0) -- (1.5,0);
				\draw (1.5,0) to [R, l^= $R$, v_=$v(t)$] (1.5,2);
			\end{circuitikz}
		\end{center}
%			\end{minipage}
%			\begin{minipage}[b]{0.33\linewidth}
		\vfill
		\begin{equation*}\label{key}
		v(t) = Ri(t)
		\end{equation*}
		\vfill
		\begin{equation*}\label{key}
		i(t) = \frac{1}{R}v(t)
		\end{equation*}
		\vfill
	\end{block}
\end{minipage}
\hfill
\begin{minipage}[b]{0.3\linewidth}
	\begin{block}{Capacitores}
		\begin{center}
			\begin{circuitikz}[european voltages]
				%\draw (0,0) to [I, l^=$i(t)$] (0,2);
				\draw [-latex] (0.6, 1.85) -- (1.25, 1.85);
				\draw node[] at (0.25, 1.75) {$i(t)$};
				\draw (0,2) -- (1.5,2);
				\draw (0,0) -- (1.5,0);
				\draw (1.5,0) to [C, l^= $C$, v_=$v(t)$] (1.5,2);
			\end{circuitikz}
		\end{center}
	%	\end{minipage}
	%	\begin{minipage}[b]{0.7\linewidth}
		\begin{equation*}\label{key}
		v(t) = \frac{1}{C}\int\limits_{ - \infty }^t {i(t)dt}
		\end{equation*}
		\begin{equation*}\label{key}
		i(t) = C\frac{dv(t)}{dt}
		\end{equation*}
		\vfill
	\end{block}
\end{minipage}
\hfill
\begin{minipage}[b]{0.3\linewidth}
	\begin{block}{Indutores}
		\begin{center}
			\begin{circuitikz}[european voltages]
				%\draw (0,0) to [I, l^=$i(t)$] (0,2);
				\draw [-latex] (0.6, 1.85) -- (1.25, 1.85);
				\draw node[] at (0.25, 1.75) {$i(t)$};
				\draw (0,2) -- (1.5,2);
				\draw (0,0) -- (1.5,0);
				\draw (1.5,0) to [L, l^= $L$, v_=$v(t)$] (1.5,2);
			\end{circuitikz}
		\end{center}
%	\end{minipage}
%	\begin{minipage}[b]{0.7\linewidth}
		\begin{equation*}\label{key}
		v(t) = L\frac{{di(t)}}{{dt}}
		\end{equation*}
		\begin{equation*}\label{key}
		i(t) = \frac{1}{L}\int\limits_{ - \infty }^t {v(t)dt}
		\end{equation*}
		\vfill
	\end{block}
\end{minipage}
\end{frame}

\begin{frame}
\frametitle{Recapitulação}
\framesubtitle{Energia armazenada no capacitor e indutor}
	\begin{minipage}[b]{0.47\linewidth}
		\begin{block}{Energia armazenada no capacitor}
			\begin{equation*}\label{key}
				\left\{ \begin{array}{l}
				p(t) = v(t)i(t) = v(t)C\frac{{dv(t)}}{{dt}}\\
				w(t) = \int\limits_{ - \infty }^t {p(t)dt} 
				\end{array} \right.
			\end{equation*}		
			Logo,
			\begin{equation*}\label{key}
				w(t) = C\int\limits_{ - \infty }^t {v(t)\frac{{dv(t)}}{{dt}}dt} 
			\end{equation*}
			\begin{equation*}\label{key}
				w(t) = C\int\limits_{ - \infty }^t {v(t)dv(t)} = \frac{1}{2}C{v^2}(t)
			\end{equation*}
		\end{block}
	\end{minipage}
	\hfill
	\begin{minipage}[b]{0.47\linewidth}
		\begin{block}{Energia armazenada no indutor}
			\begin{equation*}\label{key}
			\left\{ \begin{array}{l}
			p(t) = v(t)i(t) = i(t)L\frac{{di(t)}}{{dt}}\\
			w(t) = \int\limits_{ - \infty }^t {p(t)dt} 
			\end{array} \right.
			\end{equation*}		
			Logo,
			\begin{equation*}\label{key}
			w(t) = L\int\limits_{ - \infty }^t {i(t)\frac{{di(t)}}{{dt}}dt} 
			\end{equation*}
			\begin{equation*}\label{key}
			w(t) = L\int\limits_{ - \infty }^t {i(t)di(t)} = \frac{1}{2}L{i^2}(t)
			\end{equation*}
		\end{block}
	\end{minipage}
	
\end{frame}
%------------------------------------------------------------------------------------
\section{Funções Singulares}
	\begin{frame}
	\frametitle{Sumário}
	\small
	\tableofcontents[currentsection]
\end{frame}
%------------------------------------------------------------------------------------
\begin{frame}
\frametitle{Funções Singulares}
\begin{itemize}
	\item Utilizados para representarem sinais em circuitos com \textbf{operações chaveadas}.
	\item Auxiliam na descrição de alguns fenômenos que ocorrem na \textbf{análise de transitórios}.
	\item Por definição, são \textbf{funções descontínuas} ou que possuem derivadas descontínuas.
	\item Funções singulares mais importante: \textbf{degrau}, \textbf{rampa} e \textbf{impulso}. 
\end{itemize}
\end{frame}

	
%------------------------------------------------
\subsection{Função Degrau}
\begin{frame}
	\frametitle{Sumário}
	\small
	\tableofcontents[currentsubsection]
\end{frame}
%%-----------------------------------------------
\begin{frame}
\frametitle{Funções Singulares}
\framesubtitle{Função Degrau}
	\begin{itemize}
		\item Utilizada para representar uma variação instantânea na tensão ou na corrente;
		\vfill
		\item Possíveis aplicações: circuito de controles e sistemas digitais.
	\end{itemize}
	
\end{frame}

\begin{frame}
\frametitle{Funções Singulares}
\framesubtitle{Função Degrau}
	\begin{minipage}[b]{0.45\linewidth}
		\begin{block}{Função degrau unitária}
			\begin{equation}\label{key}
			v(t) = \left\{ \begin{array}{l}
			0,~t < 0\\
			1,~t > 0
			\end{array} \right.
			\end{equation}
			e é denotada por:
			\begin{equation}\label{key}
			v(t) = {U_{ - 1}}(t)
			\end{equation}
			\begin{center}
				\begin{circuitikz}					
					\begin{scope}[]
						\draw [-latex] (0,0) -- (3,0);
						\draw [-latex] (1.5,-0.5) -- (1.5,2);
						\draw node[] (F) at (3,-0.5) {$t$};
						\draw node[] (G) at (0.75,+2) {$U_{-1}(t)$};
						\draw [color=red] (0,0) -- (1.5,0);
						\draw [color=red] (1.5,0) -- (1.5,1);
						\draw [color=red] (1.5,1) -- (3,1);
						\draw node[] (I) at (1.25,1) {$1$};
					\end{scope}				
				\end{circuitikz}
			\end{center}
		\end{block}
	\end{minipage}
	\hfill
	\begin{minipage}[b]{0.45\linewidth}
		\begin{block}{Função degrau unitária inversa}
			\begin{equation}\label{key}
			v(t) = \left\{ \begin{array}{l}
			1,~t < 0\\
			0,~t > 0
			\end{array} \right.
			\end{equation}
			e é denotada por:
			\begin{equation}\label{key}
			v(t) = {U_{ - 1}}(-t)
			\end{equation}
			\begin{center}
				\begin{circuitikz}					
					\begin{scope}[]
						\draw [-latex] (0,0) -- (3,0);
						\draw [-latex] (1.5,-0.5) -- (1.5,2);
						\draw node[] (F) at (3,-0.5) {$t$};
						\draw node[] (G) at (0.75,+2) {$U_{-1}(-t)$};
						\draw [color=red] (0,1) -- (1.5,1);
						\draw [color=red] (1.5,1) -- (1.5,0);
						\draw [color=red] (1.5,0) -- (2.85,0);
						\draw node[] (I) at (1.75,1) {$1$};
					\end{scope}				
				\end{circuitikz}
			\end{center}
		\end{block}
	\end{minipage}
\end{frame}

\begin{frame}
\frametitle{Funções Singulares}
\framesubtitle{Função Degrau}
\newcommand{\xshift}{6}
\newcommand{\yshift}{0}
\begin{overprint}
	\only<1>
	{
		\begin{center}
			\begin{circuitikz} 
				\draw node[] at (4, 4) {para $t<t_0$};
				\draw
				(0,2) to[V, l_=$V_0$] (0,0)
				(0,2) -- (1,2)
				node[circ] (A) at (1,2) {}
				node[circ] (B) at (1.5,2) {}
				node[circ] (C) at (1.25,1.5) {}
				(1.25,1.5) -- (1.25,0)
				(0,0) -- (1.25,0)
				node[circ] (D) at (2.5,2) {}
				node[circ] (E) at (2.5,0) {}
				(B) -- (D)
				(1.25,0) -- (E)
				(2,0) -- (2.5,0)
				(D) to[open, v^=$v(t)$] (E);
				\draw 
				[-latex] (B) -- (C);					
				
				\begin{scope}[]
					\draw [-latex] (\xshift-0.5,\yshift) -- (\xshift+3,\yshift);
					\draw [-latex] (\xshift,\yshift-0.5) -- (\xshift,\yshift+2);
					\draw node[] (F) at (\xshift+3,\yshift-0.5) {$t$};
					\draw node[] (G) at (\xshift-0.5,\yshift+2) {$v(t)$};
					\draw node[] (H) at (\xshift+1.5,\yshift-0.5) {$t_0$};
					\draw [color=red] (\xshift-0.5,\yshift) -- (\xshift+1.5,\yshift);
					\draw node[] (I) at (\xshift-0.5,\yshift+1) {$V_0$};
				\end{scope}				
			\end{circuitikz}
		\end{center}
	}	
	\only<2>
	{
		\begin{center}
			\begin{circuitikz} 
				\draw node[] at (4, 4) {para $t>t_0$};
				\draw
				(0,2) to[V, l_=$V_0$] (0,0)
				(0,2) -- (1,2)
				node[circ] (A) at (1,2) {}
				node[circ] (B) at (1.5,2) {}
				node[circ] (C) at (1.25,1.5) {}
				(1.25,1.5) -- (1.25,0)
				(0,0) -- (1.25,0)
				node[circ] (D) at (2.5,2) {}
				node[circ] (E) at (2.5,0) {}
				(B) -- (D)
				(1.25,0) -- (E)
				(2,0) -- (2.5,0)
				(D) to[open, v^=$v(t)$] (E);
				\draw [-latex] (B) -- (A);					
				
				\begin{scope}[]
					\draw [-latex] (\xshift-0.5,\yshift) -- (\xshift+3,\yshift);
					\draw [-latex] (\xshift,\yshift-0.5) -- (\xshift,\yshift+2);
					\draw node[] (F) at (\xshift+3,\yshift-0.5) {$t$};
					\draw node[] (G) at (\xshift-0.5,\yshift+2) {$v(t)$};
					\draw node[] (H) at (\xshift+1.5,\yshift-0.5) {$t_0$};
					\draw [color=red] (\xshift-0.5,\yshift) -- (\xshift+1.5,\yshift);
					\draw [color=red] (\xshift+1.5,\yshift) -- (\xshift+1.5,\yshift+1);
					\draw [color=red] (\xshift+1.5,\yshift+1) -- (\xshift+3,\yshift+1);
					\draw [dotted] (\xshift,\yshift+1) -- (\xshift+1.5,\yshift+1);
					\draw node[] (I) at (\xshift-0.5,\yshift+1) {$V_0$};
				\end{scope}				
			\end{circuitikz}
		\end{center}
	}	
	\only<3>
	{
		\begin{center}
			\begin{circuitikz} 
%				\draw node[] at (4, 2.5) {para $t>t_0$};
%				\draw
%				(0,2) to[V, l_=$V_0$] (0,0)
%				(0,2) -- (1,2)
%				node[circ] (A) at (1,2) {}
%				node[circ] (B) at (1.5,2) {}
%				node[circ] (C) at (1.25,1.5) {}
%				(1.25,1.5) -- (1.25,0)
%				(0,0) -- (1.25,0)
%				node[circ] (D) at (2.5,2) {}
%				node[circ] (E) at (2.5,0) {}
%				(B) -- (D)
%				(1.25,0) -- (E)
%				(2,0) -- (2.5,0)
%				(D) to[open, v^=$v(t)$] (E);
%				\draw [-latex] (B) -- (A);	
				\begin{scope}[]
					\draw [-latex] (\xshift-0.5,\yshift) -- (\xshift+3,\yshift);
					\draw [-latex] (\xshift,\yshift-0.5) -- (\xshift,\yshift+2);
					\draw node[] (F) at (\xshift+3,\yshift-0.5) {$t$};
					\draw node[] (G) at (\xshift-0.5,\yshift+2) {$v(t)$};
					\draw node[] (H) at (\xshift+1.5,\yshift-0.5) {$t_0$};
					\draw [color=red] (\xshift-0.5,\yshift) -- (\xshift+1.5,\yshift);
					\draw [color=red] (\xshift+1.5,\yshift) -- (\xshift+1.5,\yshift+1);
					\draw [color=red] (\xshift+1.5,\yshift+1) -- (\xshift+3,\yshift+1);
					\draw [dotted] (\xshift,\yshift+1) -- (\xshift+1.5,\yshift+1);
					\draw node[] (I) at (\xshift-0.5,\yshift+1) {$V_0$};
				\end{scope}				
			\end{circuitikz}
		\end{center}
		\begin{block}{Função degrau deslocada}
			\begin{equation}\label{key}
			v(t) = \left\{ \begin{array}{l}
			0,~t < {t_0}\\
			V_0,~t > {t_0}
			\end{array} \right.
			\end{equation}
			Ou, utilizando a função degrau deslocada:
			\begin{equation}\label{key}
			v(t) = {V_0}{U_{ - 1}}(t - {t_0})
			\end{equation}
		\end{block}
	}	
\end{overprint}
\end{frame}

\begin{frame}
\frametitle{Funções Singulares}
\framesubtitle{Função Degrau}
%	\begin{block}{Função degrau deslocada}
%		\begin{equation}\label{key}
%		v(t) = \left\{ \begin{array}{l}
%		0,~t < {t_0}\\
%		V_0,~t > {t_0}
%		\end{array} \right.
%		\end{equation}
%		Ou, utilizando a função degrau deslocada:
%		\begin{equation}\label{key}
%		v(t) = {V_0}{U_{ - 1}}(t - {t_0})
%		\end{equation}
%	\end{block}
	\begin{block}{Dica!}
		Para verificar em que posição do eixo $x$ o degrau ocorre, basta igualar o valor dentro dos parenteses a zero. Exemplo, em $v(t) = 2{U_{ - 1}}(t - 5)$, deve-se fazer $t - 5 = 0 \Rightarrow t = 5$.
	\end{block}	

	\begin{center}
		\begin{circuitikz} 
			%				\draw node[] at (4, 2.5) {para $t>t_0$};
			%				\draw
			%				(0,2) to[V, l_=$V_0$] (0,0)
			%				(0,2) -- (1,2)
			%				node[circ] (A) at (1,2) {}
			%				node[circ] (B) at (1.5,2) {}
			%				node[circ] (C) at (1.25,1.5) {}
			%				(1.25,1.5) -- (1.25,0)
			%				(0,0) -- (1.25,0)
			%				node[circ] (D) at (2.5,2) {}
			%				node[circ] (E) at (2.5,0) {}
			%				(B) -- (D)
			%				(1.25,0) -- (E)
			%				(2,0) -- (2.5,0)
			%				(D) to[open, v^=$v(t)$] (E);
			%				\draw [-latex] (B) -- (A);	
			\begin{scope}[]
				\draw [-latex] (-0.5,0) -- (+3,0);
				\draw [-latex] (0,-0.5) -- (0,+2);
				\draw node[] (F) at (+3,-0.5) {$t$};
				\draw node[] (G) at (-0.5,+2) {$v(t)$};
				\draw node[] (H) at (+1.5,-0.5) {$5$};
				\draw [color=red] (-0.5,0) -- (+1.5,0);
				\draw [color=red] (+1.5,0) -- (+1.5,+1);
				\draw [color=red] (+1.5,+1) -- (+3,+1);
				\draw [dotted] (0,+1) -- (+1.5,+1);
				\draw node[] (I) at (-0.5,+1) {$2$};
			\end{scope}				
		\end{circuitikz}
	\end{center}
	

\end{frame}
%------------------------------------------------
\subsection{Função Rampa}
\begin{frame}
\frametitle{Sumário}
\small
\tableofcontents[currentsubsection]
\end{frame}
%%-----------------------------------------------
\begin{frame}
\frametitle{Funções Singulares}
\framesubtitle{Função Rampa}
	\begin{minipage}[b]{0.45\linewidth}
		\begin{block}{Função rampa unitária}
			\vspace{0.25cm}
			\begin{equation}\label{key}			
			v(t) = \left\{ \begin{array}{l}
			0,~t < 0\\
			t,~t > 0
			\end{array} \right.
			\end{equation}
			Que é denotada por:
			\begin{equation}\label{key}
			v(t) = {U_{ - 2}}(t)
			\end{equation}
			\begin{center}
				\begin{circuitikz}	
					\coordinate (Y) at (3.5,2);
					\coordinate (A) at (1.5,0);
					\coordinate (X) at (4,0);				
					\begin{scope}[]
						\draw [-latex] (0,0) -- (3,0);
						\draw [-latex] (1.5,-0.5) -- (1.5,2);
						\draw node[] (F) at (3,-0.5) {$t$};
						\draw node[] (G) at (0.75,+2) {$U_{-2}(t)$};
						\draw [color=red] (0,0) -- (1.5,0);
						\draw [color=red] (1.5,0) -- (2.5,1);
					
						\path[clip] (A) -- (X) -- (Y);
						\fill[red, opacity=0.2, draw=black] (A) circle (5mm);
						\node at ($(A)+(22.5:9mm)$) {$45^\circ$};
					\end{scope}				
				\end{circuitikz}
			\end{center}
		\end{block}
	\end{minipage}
	\hfill
	\begin{minipage}[b]{0.45\linewidth}
		\begin{block}{Função rampa genérica}
			\begin{equation}\label{key}
			v(t) = \left\{ \begin{array}{l}
			0,~t < t_0\\
			At,~t > t_0
			\end{array} \right.
			\end{equation}
			Que é denotada por:
			\begin{equation}\label{key}
			v(t) = A{U_{ - 2}}(t-t_0)
			\end{equation}
			\begin{center}
				\begin{circuitikz}	
					\coordinate (Y) at (3.5,2);
					\coordinate (A) at (1.5,0);
					\coordinate (X) at (4,0);				
					\begin{scope}[]
						\draw [-latex] (0.5,0) -- (3,0);
						\draw [-latex] (1,-0.5) -- (1,1.5);
						\draw node[] (F) at (3,-0.5) {$t$};
						\draw node[] (G) at (0.5,1.5) {$v(t)$};
						\draw [color=red] (0.5,0) -- (1.5,0);
						\draw [color=red] (1.5,0) -- (2.5,1);
						\draw node[] (F) at (1.5,-0.5) {$t_0$};
						\draw node[] (F) at (2.5,1.5) {$tg(\alpha)=A$};
						
						\path[clip] (A) -- (X) -- (Y);
						\fill[red, opacity=0.2, draw=black] (A) circle (5mm);
						\node at ($(A)+(22.5:9mm)$) {$\alpha$};
					\end{scope}				
				\end{circuitikz}
			\end{center}
		\end{block}
	\end{minipage}
\end{frame}

\begin{frame}
\frametitle{Funções Singulares}
\framesubtitle{Função Rampa}
	\begin{block}{Relação entre função rampa e função degrau}
		Note que:
		\begin{equation}\label{key}
		{U_{ - 2}}(t) = \int\limits_{ - \infty }^t {{U_{ - 1}}(t)} dt
		\end{equation}
		Ou ainda:
		\begin{equation}\label{key}
		{U_{ - 1}}(t) = \frac{{d{U_{ - 2}}(t)}}{{dt}}
		\end{equation}
	\end{block}
\end{frame}


%------------------------------------------------
\subsection{Função Impulso}
\begin{frame}
\frametitle{Sumário}
\small
\tableofcontents[currentsubsection]
\end{frame}

\begin{frame}
\frametitle{Funções Singulares}
\framesubtitle{Função Impulso}
\begin{minipage}[b]{0.45\linewidth}
	\begin{block}{Função impulso unitária}
		\vspace{0.5cm}
		\begin{equation}\label{key}			
		v(t) = \left\{ \begin{array}{l}
		0,~t \neq 0\\
		1,~t = 0
		\end{array} \right.
		\end{equation}
		Que é denotada por:
		\begin{equation}\label{key}
		v(t) = {U_0}(t)
		\end{equation}
		\begin{center}
			\begin{circuitikz}									
				\begin{scope}[]
					\draw [-latex] (0,0) -- (3,0);
					\draw [-latex] (1.5,-0.5) -- (1.5,1.5);
					\draw node[] (F) at (3,-0.5) {$t$};
					\draw node[] (G) at (0.75,1.5) {$U_{0}(t)$};
					\draw [color=red] (0,0) -- (2.85,0);
					\draw [-latex, color=red] (1.5,0) -- (1.5,1);
					\draw node[] (G) at (1.25,1) {$1$};
				\end{scope}				
			\end{circuitikz}
		\end{center}
	\end{block}
\end{minipage}
\hfill
\begin{minipage}[b]{0.45\linewidth}
	\begin{block}{Função impulso genérica}
		%		\vspace{0.25cm}
		\begin{equation}\label{key}			
		v(t) = \left\{ \begin{array}{l}
		0,~t \neq t_0\\
		A,~t = t_0
		\end{array} \right.
		\end{equation}
		Que é denotada por:
		\begin{equation}\label{key}
		v(t) = A{U_0}(t-t_0)
		\end{equation}
		\begin{center}
			\begin{circuitikz}									
				\begin{scope}[]
					\draw [-latex] (0,0) -- (3,0);
					\draw [-latex] (1,-0.5) -- (1,1.5);
					\draw node[] (F) at (3,-0.5) {$t$};
					\draw node[] (G) at (-0.25,1.5) {$AU_{0}(t-t_0)$};
					\draw [color=red] (0,0) -- (2.85,0);
					\draw [-latex, color=red] (2,0) -- (2,1);
					\draw node[] (G) at (0.75,1) {$A$};
					\draw [dotted] (1,1) -- (2,1);
					\draw node[] (F) at (2,-0.5) {$t_0$};
				\end{scope}				
			\end{circuitikz}
		\end{center}
	\end{block}
\end{minipage}
\end{frame}

\begin{frame}
	\frametitle{Funções Singulares}
	\framesubtitle{Função Rampa}
	\begin{block}{Relação entre função degrau e função impulso}
		Note que:
		\begin{equation}\label{key}
		{U_{ - 1}}(t) = \int\limits_{ - \infty }^t {{U_{0}}(t)} dt
		\end{equation}
		Ou ainda:
		\begin{equation}\label{key}
		{U_{0}}(t) = \frac{{d{U_{ - 1}}(t)}}{{dt}}
		\end{equation}
	\end{block}
\end{frame}
%%-----------------------------------------------


%------------------------------------------------
\section{Resposta às Funções Singulares}
\begin{frame}
\frametitle{Sumário}
\small
\tableofcontents[currentsection]
\end{frame}

\begin{frame}
\frametitle{Resposta às Funções Singulares}
	\begin{theorem}
		Se todas as correntes e tensões permanecerem finitas, a tensão nos terminais de uma capacitância e a corrente passando por uma indutância não podem alterar instantaneamente.
	\end{theorem}
	\begin{itemize}
		\item Tal teorema é válido para circuitos elétricos excitados pelas funções \textbf{degrau} e \textbf{rampa}.
		\item A adição de uma função \textbf{impulso} faz com que considerações adicionais sejam observadas.
	\end{itemize}
	\begin{theorem}
		Um impulso unitário de corrente entrando em uma capacitância descarregada altera a sua tensão de $\frac{1}{C}V$ e a energia de $\frac{1}{2C}J$ instantaneamente.
	\end{theorem}
\end{frame}

\begin{frame}
\frametitle{Resposta às Funções Singulares}
\newcommand{\xshift}{4}
\newcommand{\yshift}{0.3}
\begin{overprint}
	\only<1>
	{
		\begin{theorem}[Demonstração]
			\begin{center}
				\begin{circuitikz}[european voltages]
					\draw (0,0) to [I, l^=$U_0(t)~A$] (0,2);
					\draw (0,2) -- (2,2);
					\draw (0,0) -- (2,0);
					\draw (2,0) to [C, l^= $C$, v_=$v_c(t)$] (2,2);
					\begin{scope}[]
						\draw [-latex] (0+\xshift,0+\yshift) -- (3+\xshift,0+\yshift);
						\draw [-latex] (1.5+\xshift,-0.2+\yshift) -- (1.5+\xshift,1.5+\yshift);
						\draw node[] (F) at (3+\xshift,0.25+\yshift) {$t$};
						\draw node[] (G) at (0.75+\xshift,1.5+\yshift) {$U_{0}(t)$};
						\draw [color=red] (0+\xshift,0+\yshift) -- (2.85+\xshift,0+\yshift);
						\draw [-latex, color=red] (1.5+\xshift,0+\yshift) -- (1.5+\xshift,1+\yshift);
						\draw node[] (G) at (1.25+\xshift,1+\yshift) {$1$};
					\end{scope}
				\end{circuitikz}
			\end{center}
			\begin{align}\label{key}
			{v_c}(t) = \frac{1}{C}\int\limits_{ - \infty }^t {i(t)dt = } \frac{1}{C}\int\limits_{ - \infty }^{{0^ - }} {{U_0}(t)dt + } \frac{1}{C}\int\limits_{{0^ - }}^{{0^ + }} {{U_0}(t)dt + } \frac{1}{C}\int\limits_{{0^ + }}^t {{U_0}(t)dt} 
			\end{align}
		\end{theorem}
	}
	\only<2>
	{
		\begin{theorem}[Demonstração]
			\begin{center}
				\begin{circuitikz}[european voltages]
					\draw (0,0) to [I, l^=$U_0(t)~A$] (0,2);
					\draw (0,2) -- (2,2);
					\draw (0,0) -- (2,0);
					\draw (2,0) to [C, l^= $C$, v_=$v_c(t)$] (2,2);
					\begin{scope}[]
						\draw [-latex] (0+\xshift,0+\yshift) -- (3+\xshift,0+\yshift);
						\draw [-latex] (1.5+\xshift,-0.2+\yshift) -- (1.5+\xshift,1.5+\yshift);
						\draw node[] (F) at (3+\xshift,0.25+\yshift) {$t$};
						\draw node[] (G) at (0.75+\xshift,1.5+\yshift) {$U_{0}(t)$};
						\draw [color=red] (0+\xshift,0+\yshift) -- (2.85+\xshift,0+\yshift);
						\draw [-latex, color=red] (1.5+\xshift,0+\yshift) -- (1.5+\xshift,1+\yshift);
						\draw node[] (G) at (1.25+\xshift,1+\yshift) {$1$};
					\end{scope}
				\end{circuitikz}
			\end{center}			
			\begin{align}\label{key}
			{v_c}(t) = \frac{1}{C}\int\limits_{ - \infty }^t {i(t)dt = } \frac{1}{C}\int\limits_{ - \infty }^{{0^ - }} {{U_0}(t)dt + } \frac{1}{C}\int\limits_{{0^ - }}^{{0^ + }} {{U_0}(t)dt + } \textcolor{red}{\frac{1}{C}\int\limits_{{0^ + }}^t {{U_0}(t)dt}} 
			\end{align}
			\textcolor{red}{Componente nula, já que o impulso ocorreu anteriormente.}
		\end{theorem}	
	}
	\only<3>
	{
		\begin{theorem}[Demonstração]
			\begin{center}
				\begin{circuitikz}[european voltages]
					\draw (0,0) to [I, l^=$U_0(t)~A$] (0,2);
					\draw (0,2) -- (2,2);
					\draw (0,0) -- (2,0);
					\draw (2,0) to [C, l^= $C$, v_=$v_c(t)$] (2,2);
					\begin{scope}[]
						\draw [-latex] (0+\xshift,0+\yshift) -- (3+\xshift,0+\yshift);
						\draw [-latex] (1.5+\xshift,-0.2+\yshift) -- (1.5+\xshift,1.5+\yshift);
						\draw node[] (F) at (3+\xshift,0.25+\yshift) {$t$};
						\draw node[] (G) at (0.75+\xshift,1.5+\yshift) {$U_{0}(t)$};
						\draw [color=red] (0+\xshift,0+\yshift) -- (2.85+\xshift,0+\yshift);
						\draw [-latex, color=red] (1.5+\xshift,0+\yshift) -- (1.5+\xshift,1+\yshift);
						\draw node[] (G) at (1.25+\xshift,1+\yshift) {$1$};
					\end{scope}
				\end{circuitikz}
			\end{center}			
			\begin{align}\label{key}
			{v_c}(t) = \frac{1}{C}\int\limits_{ - \infty }^t {i(t)dt = } \textcolor{red}{\frac{1}{C}\int\limits_{ - \infty }^{{0^ - }}
			{{U_0}(t)dt}} + \frac{1}{C}\int\limits_{{0^ - }}^{{0^ + }} {{U_0}(t)dt + }  
			0
			\end{align}
			\textcolor{red}{Componente que representa a energia inicial armazenada do capacitor. Componente nula se o capacitor for considerado inicialmente descarregado.}
		\end{theorem}	
	}
	\only<4>
	{
		\begin{theorem}[Demonstração]
			\begin{center}
				\begin{circuitikz}[european voltages]
					\draw (0,0) to [I, l^=$U_0(t)~A$] (0,2);
					\draw (0,2) -- (2,2);
					\draw (0,0) -- (2,0);
					\draw (2,0) to [C, l^= $C$, v_=$v_c(t)$] (2,2);
					\begin{scope}[]
						\draw [-latex] (0+\xshift,0+\yshift) -- (3+\xshift,0+\yshift);
						\draw [-latex] (1.5+\xshift,-0.2+\yshift) -- (1.5+\xshift,1.5+\yshift);
						\draw node[] (F) at (3+\xshift,0.25+\yshift) {$t$};
						\draw node[] (G) at (0.75+\xshift,1.5+\yshift) {$U_{0}(t)$};
						\draw [color=red] (0+\xshift,0+\yshift) -- (2.85+\xshift,0+\yshift);
						\draw [-latex, color=red] (1.5+\xshift,0+\yshift) -- (1.5+\xshift,1+\yshift);
						\draw node[] (G) at (1.25+\xshift,1+\yshift) {$1$};
					\end{scope}
				\end{circuitikz}
			\end{center}			
			\begin{equation}\label{key}
			{v_c}(t) = \frac{1}{C}\int\limits_{ - \infty }^t {i(t)dt = } 0 + \frac{1}{C}\int\limits_{{0^ - }}^{{0^ + }} {{U_0}(t)dt + } 0 = \frac{1}{C}{U_{ - 1}}(t) = 
			\textcolor{red}{\frac{1}{C}~V}
			\end{equation}
			e
			\begin{equation}\label{key}
			w(t) = \frac{1}{2}Cv_c^2(t) = \frac{1}{2}C{\left( {\frac{1}{C}{U_{ - 1}}(t)} \right)^2} = \frac{1}{{2C}}{U_{ - 1}}(t)
			= \textcolor{red}{\frac{1}{{2C}}~J}
			\end{equation}
			\textcolor{red}{Para $t>0$.}
		\end{theorem}			
	}
\end{overprint}
\end{frame}

%------------------------------------------------
\section{Decomposição de Sinais em Funções Singulares}
\begin{frame}
\frametitle{Sumário}
\small
\tableofcontents[currentsection]
\end{frame}

\begin{frame}
\frametitle{Decomposição de Sinais em Funções Singulares}
	\begin{itemize}
		\item Uma forma de onda de tensão ou corrente, constituída por segmentos de reta, pode ser \textbf{decomposta em finitas funções degraus ou rampa}.
		\vfill
		\item Utilização do \textbf{teorema da superposição}.
	\end{itemize}	
\end{frame}

\begin{frame}
\frametitle{Decomposição de Sinais em Funções Singulares}
\begin{overprint}
	\only<1>
	{
		\begin{block}{Exemplo 1}
			\begin{center}
				\begin{circuitikz} 			
					\begin{scope}[]
						\draw [-latex] (-0.5,0) -- (2,0);
						\draw [-latex] (0,-0.25) -- (0,2);
						\draw node[] (F) at (2,-0.25) {$t$};
						\draw node[] (G) at (-0.5,1.75) {$x(t)$};
						\draw node[] (H) at (+1,-0.25) {$t_0$};
						\draw [color=red] (-0.5,0) -- (0,0);
						\draw [color=red] (0,0) -- (0,1);
						\draw [color=red] (0,1) -- (1,1);
						\draw [color=red] (1,1) -- (1,0);
						\draw [color=red] (1,0) -- (1.85,0);
						\draw node[] (I) at (-0.5,+1) {$X_m$};
					\end{scope}				
				\end{circuitikz}
			\end{center}		
		\end{block}
	}
	\only<2>
	{
		\begin{block}{Exemplo 1}
			\begin{center}
				\begin{circuitikz}[scale=0.5, every node/.style={scale=0.5}]		
					\begin{scope}[]
						\draw [-latex] (-0.5,0) -- (2,0);
						\draw [-latex] (0,-0.25) -- (0,2);
						\draw node[] (F) at (2,-0.25) {$t$};
						\draw node[] (G) at (-0.5,1.75) {$x(t)$};
						\draw node[] (H) at (+1,-0.25) {$t_0$};
						\draw [color=red] (-0.5,0) -- (0,0);
						\draw [color=red] (0,0) -- (0,1);
						\draw [color=red] (0,1) -- (1,1);
						\draw [color=red] (1,1) -- (1,0);
						\draw [color=red] (1,0) -- (1.85,0);
						\draw node[] (I) at (-0.5,+1) {$X_m$};
					\end{scope}				
				\end{circuitikz}
			\end{center}
			\begin{minipage}[b]{0.4\linewidth}
				\begin{center}
					\begin{circuitikz} 			
						\begin{scope}[]
							\draw [-latex] (-0.5,0) -- (2,0);
							\draw [-latex] (0,-0.25) -- (0,2);
							\draw node[] (F) at (2,-0.25) {$t$};
							\draw node[] (G) at (-0.5,1.75) {$x(t)$};
							\draw node[] (H) at (+1,-0.25) {$t_0$};
							\draw [color=purple] (-0.5,0) -- (0,0);
							\draw [color=purple] (0,0) -- (0,1);
							\draw [color=purple] (0,1) -- (1.5,1);
							\draw [color=purple, dotted] (1.5,1) -- (2,1);
							\draw node[] (I) at (-0.5,+1) {$X_m$};
						\end{scope}				
					\end{circuitikz}
				\end{center}
			\end{minipage}	
			\hfill	
			\begin{minipage}[b]{0.1\linewidth}
				\vspace{1cm}
				\begin{center}
					$\Rightarrow$
				\end{center}
				\vspace{1cm}				
			\end{minipage}	
			\hfill
			\begin{minipage}[b]{0.4\linewidth}
				\begin{center}
					\begin{circuitikz} 			
						\begin{scope}[]
							\draw [-latex] (-0.5,0) -- (2,0);
							\draw [-latex] (0,-0.25) -- (0,2);
							\draw node[] (F) at (2,-0.25) {$t$};
							\draw node[] (G) at (-0.5,1.75) {$x(t)$};
							\draw node[] (H) at (+1,-0.25) {$t_0$};
							\draw [color=red] (-0.5,0) -- (0,0);
							\draw [color=red] (0,0) -- (0,1);
							\draw [color=red] (0,1) -- (1.5,1);
							\draw [color=red, dotted] (1.5,1) -- (2,1);
							\draw node[] (I) at (-0.5,+1) {$X_m$};
						\end{scope}			
					\end{circuitikz}
				\end{center}
			\end{minipage}	
			\begin{center}
				$x(t)=\textcolor{purple}{X_mU_{-1}(t)}$
			\end{center}
		\end{block}		
	}
	\only<3>
	{
		\begin{block}{Exemplo 1}
			\begin{center}
				\begin{circuitikz}[scale=0.5, every node/.style={scale=0.5}]		
					\begin{scope}[]
						\draw [-latex] (-0.5,0) -- (2,0);
						\draw [-latex] (0,-0.25) -- (0,2);
						\draw node[] (F) at (2,-0.25) {$t$};
						\draw node[] (G) at (-0.5,1.75) {$x(t)$};
						\draw node[] (H) at (+1,-0.25) {$t_0$};
						\draw [color=red] (-0.5,0) -- (0,0);
						\draw [color=red] (0,0) -- (0,1);
						\draw [color=red] (0,1) -- (1,1);
						\draw [color=red] (1,1) -- (1,0);
						\draw [color=red] (1,0) -- (1.85,0);
						\draw node[] (I) at (-0.5,+1) {$X_m$};
					\end{scope}				
				\end{circuitikz}
			\end{center}
			\begin{minipage}[b]{0.4\linewidth}
				\begin{center}
					\begin{circuitikz} 			
						\begin{scope}[]
							\draw [-latex] (-0.5,0) -- (2,0);
							\draw [-latex] (0,-1.25) -- (0,2);
							\draw node[] (F) at (2,-0.25) {$t$};
							\draw node[] (G) at (-0.5,1.75) {$x(t)$};
							\draw node[] (H) at (+0.75,-0.25) {$t_0$};
							\draw [color=purple] (-0.5,0) -- (0,0);
							\draw [color=purple] (0,0) -- (0,1);
							\draw [color=purple] (0,1) -- (1.5,1);
							\draw [color=purple, dotted] (1.5,1) -- (2,1);
							\draw [color=blue] (0,0) -- (1,0);
							\draw [color=blue] (1,0) -- (1,-1);
							\draw [color=blue] (1,-1) -- (1.5,-1); 
							\draw [color=blue, dotted] (1.5,-1) -- (2,-1);
							\draw node[] (I) at (-0.5,+1) {$X_m$};
							\draw node[] (I) at (-0.5,-1) {$-X_m$};
						\end{scope}				
					\end{circuitikz}
				\end{center}
			\end{minipage}	
			\hfill	
			\begin{minipage}[b]{0.1\linewidth}
				\vspace{1cm}
				\begin{center}
					$\Rightarrow$
				\end{center}
				\vspace{1cm}				
			\end{minipage}	
			\hfill
			\begin{minipage}[b]{0.4\linewidth}
				\begin{center}
					\begin{circuitikz} 			
						\begin{scope}[]
							\draw [-latex] (-0.5,0) -- (2,0);
							\draw [-latex] (0,-0.25) -- (0,2);
							\draw node[] (F) at (2,-0.25) {$t$};
							\draw node[] (G) at (-0.5,1.75) {$x(t)$};
							\draw node[] (H) at (+1,-0.25) {$t_0$};
							\draw [color=red] (-0.5,0) -- (0,0);
							\draw [color=red] (0,0) -- (0,1);
							\draw [color=red] (0,1) -- (1,1);
							\draw [color=red] (1,1) -- (1,0);
							\draw [color=red] (1,0) -- (1.85,0);
							\draw node[] (I) at (-0.5,+1) {$X_m$};
						\end{scope}			
					\end{circuitikz}
				\end{center}
			\end{minipage}	
			\begin{center}
				$x(t)=\textcolor{purple}{X_mU_{-1}(t)} \textcolor{blue}{- X_mU_{-1}(t-t_0)}$ 
			\end{center}
		\end{block}		
	}
\end{overprint}
\end{frame}

\begin{frame}
\frametitle{Decomposição de Sinais em Funções Singulares}
\begin{overprint}
	\only<1>
	{
		\begin{block}{Exemplo 2}
			\begin{center}
				\begin{circuitikz} 			
					\begin{scope}[]
						\draw [-latex] (-0.5,0) -- (2,0);
						\draw [-latex] (0,-0.25) -- (0,2);
						\draw node[] (F) at (2,-0.25) {$t$};
						\draw node[] (G) at (-0.5,1.75) {$x(t)$};
						\draw node[] (H) at (+1,-0.25) {$2$};
						\draw node[] (H) at (+1.5,-0.25) {$3$};
						\draw node[] (I) at (-0.25,0.5) {$1$};
						\draw node[] (I) at (-0.25,1) {$2$};
						\draw [color=red] (-0.5,0) -- (0,0);
						\draw [color=red] (0,0) -- (1,1);
						\draw [color=red] (1,1) -- (1,0.5);
						\draw [color=red] (1,0.5) -- (1.5,0.5);
						\draw [color=red] (1.5,0.5) -- (1.5,0);
						\draw [color=red] (1.5,0) -- (1.85,0);
						
					\end{scope}				
				\end{circuitikz}
			\end{center}
		\end{block}
	}
	\only<2>
	{
		\begin{block}{Exemplo 2}
			\begin{center}
				\begin{circuitikz}[scale=0.5, every node/.style={scale=0.5}]		
					\begin{scope}[]
						\draw [-latex] (-0.5,0) -- (2,0);
						\draw [-latex] (0,-0.25) -- (0,2);
						\draw node[] (F) at (2,-0.25) {$t$};
						\draw node[] (G) at (-0.5,1.75) {$x(t)$};
						\draw node[] (H) at (+1,-0.25) {$2$};
						\draw node[] (H) at (+1.5,-0.25) {$3$};
						\draw node[] (I) at (-0.25,0.5) {$1$};
						\draw node[] (I) at (-0.25,1) {$2$};
						\draw [color=red] (-0.5,0) -- (0,0);
						\draw [color=red] (0,0) -- (1,1);
						\draw [color=red] (1,1) -- (1,0.5);
						\draw [color=red] (1,0.5) -- (1.5,0.5);
						\draw [color=red] (1.5,0.5) -- (1.5,0);
						\draw [color=red] (1.5,0) -- (1.7,0);
					\end{scope}				
				\end{circuitikz}
			\end{center}
			\begin{minipage}[b]{0.4\linewidth}
				\begin{center}
					\begin{circuitikz} 			
						\begin{scope}[]
							\draw [-latex] (-0.5,0) -- (3,0);
							\draw [-latex] (0,-0.25) -- (0,2);
							\draw node[] (F) at (3,-0.25) {$t$};
							\draw node[] (G) at (-0.5,1.75) {$x(t)$};
							\draw node[] (H) at (+1,-0.25) {$2$};
							\draw node[] (H) at (+1.5,-0.25) {$3$};
							\draw node[] (I) at (-0.25,0.5) {$1$};
							\draw node[] (I) at (-0.25,1) {$2$};
							\draw [color=purple] (-0.5,0) -- (0,0);
							\draw [color=purple] (0,0) -- (1,1);
							\draw [color=purple, dotted] (1,1) -- (1.5,1.5);
						\end{scope}				
					\end{circuitikz}
				\end{center}
			\end{minipage}	
			\hfill	
			\begin{minipage}[b]{0.1\linewidth}
				\vspace{1cm}
				\begin{center}
					$\Rightarrow$
				\end{center}
				\vspace{1cm}				
			\end{minipage}	
			\hfill
			\begin{minipage}[b]{0.4\linewidth}
				\begin{center}
					\begin{circuitikz} 			
						\begin{scope}[]
							\draw [-latex] (-0.5,0) -- (2,0);
							\draw [-latex] (0,-0.25) -- (0,2);
							\draw node[] (F) at (2,-0.25) {$t$};
							\draw node[] (G) at (-0.5,1.75) {$x(t)$};
							\draw node[] (H) at (+1,-0.25) {$2$};
							\draw node[] (H) at (+1.5,-0.25) {$3$};
							\draw node[] (I) at (-0.25,0.5) {$1$};
							\draw node[] (I) at (-0.25,1) {$2$};
							\draw [color=red] (-0.5,0) -- (0,0);
							\draw [color=red] (0,0) -- (1,1);
							\draw [color=red, dotted] (1,1) -- (1.5,1.5);
						\end{scope}			
					\end{circuitikz}
				\end{center}
			\end{minipage}	
			\begin{center}
				$x(t)=\textcolor{purple}{U_{-2}(t)}$ 
			\end{center}
		\end{block}	
	}
	\only<3>
	{
		\begin{block}{Exemplo 2}
			\begin{center}
				\begin{circuitikz}[scale=0.5, every node/.style={scale=0.5}]		
					\begin{scope}[]
						\draw [-latex] (-0.5,0) -- (2,0);
						\draw [-latex] (0,-0.25) -- (0,2);
						\draw node[] (F) at (2,-0.25) {$t$};
						\draw node[] (G) at (-0.5,1.75) {$x(t)$};
						\draw node[] (H) at (+1,-0.25) {$2$};
						\draw node[] (H) at (+1.5,-0.25) {$3$};
						\draw node[] (I) at (-0.25,0.5) {$1$};
						\draw node[] (I) at (-0.25,1) {$2$};
						\draw [color=red] (-0.5,0) -- (0,0);
						\draw [color=red] (0,0) -- (1,1);
						\draw [color=red] (1,1) -- (1,0.5);
						\draw [color=red] (1,0.5) -- (1.5,0.5);
						\draw [color=red] (1.5,0.5) -- (1.5,0);
						\draw [color=red] (1.5,0) -- (1.7,0);
					\end{scope}				
				\end{circuitikz}
			\end{center}
			\begin{minipage}[b]{0.4\linewidth}
				\begin{center}
					\begin{circuitikz} 			
						\begin{scope}[]
							\draw [-latex] (-0.5,0) -- (3,0);
							\draw [-latex] (0,-0.25) -- (0,2);
							\draw node[] (F) at (3,0.25) {$t$};
							\draw node[] (G) at (-0.5,1.75) {$x(t)$};
							\draw node[] (H) at (+1,0.25) {$2$};
							\draw node[] (H) at (+1.5,0.25) {$3$};
							\draw node[] (I) at (-0.25,0.5) {$1$};
							\draw node[] (I) at (-0.25,1) {$2$};
							\draw [color=purple] (-0.5,0) -- (0,0);
							\draw [color=purple] (0,0) -- (1,1);
							\draw [color=purple, dotted] (1,1) -- (1.5,1.5);
							\draw [color=blue] (0,0) -- (1,0);
							\draw [color=blue] (1,0) -- (1.75,-0.75);
							\draw [color=blue, dotted] (1.75,-0.75) -- (2,-1);
						\end{scope}				
					\end{circuitikz}
				\end{center}
			\end{minipage}	
			\hfill	
			\begin{minipage}[b]{0.1\linewidth}
				\vspace{1cm}
				\begin{center}
					$\Rightarrow$
				\end{center}
				\vspace{1cm}				
			\end{minipage}	
			\hfill
			\begin{minipage}[b]{0.4\linewidth}
				\begin{center}
					\begin{circuitikz} 			
						\begin{scope}[]
							\draw [-latex] (-0.5,0) -- (2,0);
							\draw [-latex] (0,-0.25) -- (0,2);
							\draw node[] (F) at (2,-0.25) {$t$};
							\draw node[] (G) at (-0.5,1.75) {$x(t)$};
							\draw node[] (H) at (+1,-0.25) {$2$};
							\draw node[] (H) at (+1.5,-0.25) {$3$};
							\draw node[] (I) at (-0.25,0.5) {$1$};
							\draw node[] (I) at (-0.25,1) {$2$};
							\draw [color=red] (-0.5,0) -- (0,0);
							\draw [color=red] (0,0) -- (1,1);
							\draw [color=red, dotted] (1,1) -- (2,1);
						\end{scope}			
					\end{circuitikz}
				\end{center}
			\end{minipage}	
			\begin{center}
				$x(t)=\textcolor{purple}{U_{-2}(t)} \textcolor{blue}{- U_{-2}(t-2)}$
			\end{center}
		\end{block}	
	}
	\only<4>
	{
		\begin{block}{Exemplo 2}
			\begin{center}
				\begin{circuitikz}[scale=0.5, every node/.style={scale=0.5}]		
					\begin{scope}[]
						\draw [-latex] (-0.5,0) -- (2,0);
						\draw [-latex] (0,-0.25) -- (0,2);
						\draw node[] (F) at (2,-0.25) {$t$};
						\draw node[] (G) at (-0.5,1.75) {$x(t)$};
						\draw node[] (H) at (+1,-0.25) {$2$};
						\draw node[] (H) at (+1.5,-0.25) {$3$};
						\draw node[] (I) at (-0.25,0.5) {$1$};
						\draw node[] (I) at (-0.25,1) {$2$};
						\draw [color=red] (-0.5,0) -- (0,0);
						\draw [color=red] (0,0) -- (1,1);
						\draw [color=red] (1,1) -- (1,0.5);
						\draw [color=red] (1,0.5) -- (1.5,0.5);
						\draw [color=red] (1.5,0.5) -- (1.5,0);
						\draw [color=red] (1.5,0) -- (1.7,0);
					\end{scope}				
				\end{circuitikz}
			\end{center}
			\begin{minipage}[b]{0.4\linewidth}
				\begin{center}
					\begin{circuitikz} 			
						\begin{scope}[]
							\draw [-latex] (-0.5,0) -- (3,0);
							\draw [-latex] (0,-0.25) -- (0,2);
							\draw node[] (F) at (3,0.25) {$t$};
							\draw node[] (G) at (-0.5,1.75) {$x(t)$};
							\draw node[] (H) at (+1,0.25) {$2$};
							\draw node[] (H) at (+1.5,0.25) {$3$};
							\draw node[] (I) at (-0.25,0.5) {$1$};
							\draw node[] (I) at (-0.25,1) {$2$};
							\draw [color=purple] (-0.5,0) -- (0,0);
							\draw [color=purple] (0,0) -- (1,1);
							\draw [color=purple, dotted] (1,1) -- (1.5,1.5);
							\draw [color=blue] (0,0) -- (1,0);
							\draw [color=blue] (1,0) -- (1.75,-0.75);
							\draw [color=blue, dotted] (1.75,-0.75) -- (2,-1);
							\draw [color=Green] (0,-0.01) -- (1,-0.01);
							\draw [color=Green] (1,-0.01) -- (1,-0.5);
							\draw [color=Green] (1,-0.5) -- (2.5,-0.5);
							\draw [color=Green, dotted] (2.5,-0.5) -- (3,-0.5);
						\end{scope}				
					\end{circuitikz}
				\end{center}
			\end{minipage}	
			\hfill	
			\begin{minipage}[b]{0.1\linewidth}
				\vspace{1cm}
				\begin{center}
					$\Rightarrow$
				\end{center}
				\vspace{1cm}				
			\end{minipage}	
			\hfill
			\begin{minipage}[b]{0.4\linewidth}
				\begin{center}
					\begin{circuitikz} 			
						\begin{scope}[]
							\draw [-latex] (-0.5,0) -- (2,0);
							\draw [-latex] (0,-0.25) -- (0,2);
							\draw node[] (F) at (2,-0.25) {$t$};
							\draw node[] (G) at (-0.5,1.75) {$x(t)$};
							\draw node[] (H) at (+1,-0.25) {$2$};
							\draw node[] (H) at (+1.5,-0.25) {$3$};
							\draw node[] (I) at (-0.25,0.5) {$1$};
							\draw node[] (I) at (-0.25,1) {$2$};
							\draw [color=red] (-0.5,0) -- (0,0);
							\draw [color=red] (0,0) -- (1,1);
							\draw [color=red] (1,1) -- (1,0.5);
							\draw [color=red] (1,0.5) -- (1.5,0.5);
							\draw [color=red, dotted] (1.5,0.5) -- (2,0.5);
						\end{scope}			
					\end{circuitikz}
				\end{center}
			\end{minipage}	
			\begin{center}
				$x(t)=\textcolor{purple}{U_{-2}(t)} \textcolor{blue}{- U_{-2}(t-2)} \textcolor{Green}{- U_{-1}(t-2)}$
			\end{center}
		\end{block}	
	}
	\only<5>
	{
		\begin{block}{Exemplo 2}
			\begin{center}
				\begin{circuitikz}[scale=0.5, every node/.style={scale=0.5}]		
					\begin{scope}[]
						\draw [-latex] (-0.5,0) -- (2,0);
						\draw [-latex] (0,-0.25) -- (0,2);
						\draw node[] (F) at (2,-0.25) {$t$};
						\draw node[] (G) at (-0.5,1.75) {$x(t)$};
						\draw node[] (H) at (+1,-0.25) {$2$};
						\draw node[] (H) at (+1.5,-0.25) {$3$};
						\draw node[] (I) at (-0.25,0.5) {$1$};
						\draw node[] (I) at (-0.25,1) {$2$};
						\draw [color=red] (-0.5,0) -- (0,0);
						\draw [color=red] (0,0) -- (1,1);
						\draw [color=red] (1,1) -- (1,0.5);
						\draw [color=red] (1,0.5) -- (1.5,0.5);
						\draw [color=red] (1.5,0.5) -- (1.5,0);
						\draw [color=red] (1.5,0) -- (1.7,0);
					\end{scope}				
				\end{circuitikz}
			\end{center}
			\begin{minipage}[b]{0.4\linewidth}
				\begin{center}
					\begin{circuitikz} 			
						\begin{scope}[]
							\draw [-latex] (-0.5,0) -- (3,0);
							\draw [-latex] (0,-0.25) -- (0,2);
							\draw node[] (F) at (3,0.25) {$t$};
							\draw node[] (G) at (-0.5,1.75) {$x(t)$};
							\draw node[] (H) at (+1,0.25) {$2$};
							\draw node[] (H) at (+1.5,0.25) {$3$};
							\draw node[] (I) at (-0.25,0.5) {$1$};
							\draw node[] (I) at (-0.25,1) {$2$};
							\draw [color=purple] (-0.5,0) -- (0,0);
							\draw [color=purple] (0,0) -- (1,1);
							\draw [color=purple, dotted] (1,1) -- (1.5,1.5);
							\draw [color=blue] (0,0) -- (1,0);
							\draw [color=blue] (1,0) -- (1.75,-0.75);
							\draw [color=blue, dotted] (1.75,-0.75) -- (2,-1);
							\draw [color=Green] (0,-0.01) -- (1,-0.01);
							\draw [color=Green] (1,-0.01) -- (1,-0.5);
							\draw [color=Green] (1,-0.5) -- (2.5,-0.5);
							\draw [color=Green, dotted] (2.5,-0.5) -- (3,-0.5);
							\draw [color=YellowOrange] (0,-0.02) -- (1.5,-0.02);
							\draw [color=YellowOrange] (1.5,-0.02) -- (1.5,-0.51);
							\draw [color=YellowOrange] (1.5,-0.51) -- (2.5,-0.51);
							\draw [color=YellowOrange, dotted] (2.5,-0.51) -- (3,-0.51);
						\end{scope}				
					\end{circuitikz}
				\end{center}
			\end{minipage}	
			\hfill	
			\begin{minipage}[b]{0.1\linewidth}
				\vspace{1cm}
				\begin{center}
					$\Rightarrow$
				\end{center}
				\vspace{1cm}				
			\end{minipage}	
			\hfill
			\begin{minipage}[b]{0.4\linewidth}
				\begin{center}
					\begin{circuitikz} 			
						\begin{scope}[]
							\draw [-latex] (-0.5,0) -- (2,0);
							\draw [-latex] (0,-0.25) -- (0,2);
							\draw node[] (F) at (2,-0.25) {$t$};
							\draw node[] (G) at (-0.5,1.75) {$x(t)$};
							\draw node[] (H) at (+1,-0.25) {$2$};
							\draw node[] (H) at (+1.5,-0.25) {$3$};
							\draw node[] (I) at (-0.25,0.5) {$1$};
							\draw node[] (I) at (-0.25,1) {$2$};
							\draw [color=red] (-0.5,0) -- (0,0);
							\draw [color=red] (0,0) -- (1,1);
							\draw [color=red] (1,1) -- (1,0.5);
							\draw [color=red] (1,0.5) -- (1.5,0.5);
							\draw [color=red] (1.5,0.5) -- (1.5,0);
							\draw [color=red] (1.5,0) -- (1.85,0);
						\end{scope}			
					\end{circuitikz}
				\end{center}
			\end{minipage}	
			\begin{center}
				$x(t)=\textcolor{purple}{U_{-2}(t)} \textcolor{blue}{- U_{-2}(t-2)} \textcolor{Green}{- U_{-1}(t-2)} \textcolor{YellowOrange}{- U_{-1}(t-3)}$
			\end{center}
		\end{block}	
	}
\end{overprint}
\end{frame}

%------------------------------------------------
\begin{frame}
\frametitle{References}
\footnotesize{
	\begin{thebibliography}{99} % Beamer does not support BibTeX so references must be inserted manually as below
		\bibitem[Vander Menengoy da Costa, 2013]{p1} Vander Menengoy da Costa (2013).
		\newblock Circuitos elétricos lineares: enfoque teórico e prático.
		\newblock \emph{Editora Interciência}.
		
		\bibitem[Charles M. Close, 1975]{p1} Charles M. Close (1975).
		\newblock Circuitos Lineares.
		\newblock \emph{Livros Técnicos e Científicos Editora S.A.}.
	\end{thebibliography}
}
\end{frame}

\end{document}
